% base/package

%
% Codificación
%
% Codificación de salida: acentos, ñ, etc. como glifos reales
\usepackage[T1]{fontenc}
% Codificación de entrada: acepta ó, ñ directamente en el .tex
% En opciones modernas es por default
\usepackage[utf8]{inputenc}       


%
% Idioma
%
% Tipografía según idioma activo (guiones, fechas, títulos)
\usepackage[english, spanish]{babel}  
% Comillas tipográficas correctas; requerido por biblatex con babel
\usepackage{csquotes}             


%
% Utilidades
%
% Colores: \textcolor{red}{...}, \colorbox{}, definición de colores propios
\usepackage{xcolor}     
% Herramientas para programación en LaTeX: \AtBeginEnvironment, \patchcmd, etc.
\usepackage{etoolbox}   
% Lorem ipsum
\usepackage{lipsum}
% Para comentarios multilinea
\usepackage{comment}
% Para inserción de pdf
\usepackage{pdfpages}


%
% Matemáticas
%
% Fuentes matemáticas: ℝ, ℕ, ℤ con \mathbb{} (carga amstext)
\usepackage{amsmath}   
% Símbolos extra: ∅, ∈, ⊆, ≤, etc. (carga amsfonts)
\usepackage{amssymb}   
% Entornos: theorem, proof, lemma, definition, etc.
\usepackage{amsthm}    
% Extiende amsmath: \coloneqq, corchetes extensibles, etc.
\usepackage{mathtools} 
% Diagramas conmutativos simples (flechas →, ↓)
\usepackage{amscd}     
% Define símbolos matemáticos en negrita de forma elegante
\usepackage{bm}


%
% Fuente
%
% Para Times New Roman (texto y partes matemáticas)
\usepackage{newtxtext, newtxmath}
% Para tipografía fina
\usepackage{microtype}


%
% Layout y geometría
%
% Control total del tamaño de página y márgenes
\usepackage{geometry}   
% Herramientas para manejar interlineado
\usepackage{setspace}
% Encabezados y pies de página personalizados
\usepackage{fancyhdr}   
% Rediseña el formato de \chapter, \section, etc.
\usepackage{titlesec}
% Redefine el formato de los títulos en el índice (toc)
\usepackage{titletoc}
% Provee \pageref{LastPage} para "Página X de Y"
\usepackage{lastpage}   


%
% Figuras, tablas y entornos
%
% \includegraphics{} con opciones (scale, width, angle)
\usepackage{graphicx}   
% Especificador [H]: coloca el flotante exactamente aquí
\usepackage{float}      
% Texto que rodea una figura lateralmente
\usepackage{wrapfig}    
% Personaliza el estilo de los \caption{}
\usepackage{caption}    
% Cajas con borde que pueden partirse entre páginas
\usepackage{framed}     
% Para tablas largas que se parten entre páginas
\usepackage{longtable}
% Para \newcolumntype
\usepackage{array}
% \RaggedRight con shrink finito (evita "Infinite glue shrinkage" al partir
% entre páginas tablas largas con columnas en bandera)
\usepackage{ragged2e}
% Para multirow
\usepackage{multirow}
% Páginas en orientación horizontal (landscape), para tablas/figuras anchas
\usepackage{pdflscape}
% Diagramas de Gantt para cronogramas (primer paquete que trae TikZ/PGF a tu preámbulo)
\usepackage{pgfgantt}


%
% Listas
%
% Control fino de listas: itemize, enumerate, description
\usepackage{enumitem}   


%
% Fuentes
%
% Letras caligráficas tipo Euler (\EuScript{A})
\usepackage{euscript}


%
% Referencias e índice
%
% Índice analítico (admite múltiples índices)
\usepackage{imakeidx}                                   
% Links en PDF: refs, URLs, TOC clicables
\usepackage{hyperref}                                   
% \cref{} que escribe "ecuación (3)" automáticamente
\usepackage[nameinlink]{cleveref}                       
% Bibliografía moderna; biber como motor
\usepackage[backend=biber, style=apa, language=spanish, sorting=nyt]{biblatex}


